\documentclass{article}
\usepackage{graphicx}

\usepackage[utf8]{inputenc} % Para aceitar acentos
\usepackage[T1]{fontenc}    % Para renderizar fontes com acentos corretamente
\usepackage[brazil]{babel}  % Para português (hifenização, datas, etc.)
\usepackage[a4paper, margin=2.5cm]{geometry} % Ajusta a largura da folha
\usepackage[ruled,vlined, linesnumbered]{algorithm2e}% Escrever pseudocódigo
\usepackage{amsmath}
\usepackage{amssymb}
\usepackage{hyperref}
\usepackage{xcolor}
\newcommand{\blue}[1]{\textcolor{blue}{#1}}
\newcommand{\red}[1]{\textcolor{red}{#1}}

\setlength{\parindent}{0pt}
\setlength{\parskip}{4pt}

\title{Lista 03 - Introdução a Criptografia}
\author{Antônio Erick Freitas Ferreira - 322980\\
        \small Instituto de Computação - Unicamp - São Paulo - Brasil
}

\begin{document}
\maketitle

\textbf{Questão 1}

    Seja $k$ a chave descoberta e $c_1$ e $c_2$ os arquivos capturados na rede. Como $c_2 = ENC_k (0xFF \oplus IV)$, faça $y_2 = DEC_k(c_2) = 0xFF \oplus IV$. Agora extraía $IV$ da seguinte forma: $IV = y_2 \oplus 0xFF = 0xFF \oplus IV \oplus 0xFF = IV$.

    Para recuperar o conteúdo da mensagem $c_1$, temos que $c_1 = ENC_k(m_1 \oplus IV)$, faça $y_1 = DEC_k(c_1) = m_1 \oplus IV$ e depois $m_1 = y_1 \oplus IV = m_1 \oplus IV \oplus IV = m_1$.

\textbf{Questão 2}

\textbf{Item A)}

    \begin{algorithm}[H]
        \caption{Desencriptador}
        \KwIn{$(IV, c_1, \cdots, c_l)$}
        \KwOut{$m \in \{0,1\}^{nl}$}
    
        $L \leftarrow []$

        \For{$1 \leq i \leq l$}{
            $y \leftarrow F_k^{-1}(c_i)$
            
            $x \leftarrow (y - IV - i)$ mod $2^n$  

            Insira $x$ em $L$
        }

        $m \leftarrow L_1||L_2||\cdots||L_l$

        \Return $m$
    \end{algorithm}


\textbf{Item B)}

    \begin{algorithm}[H]
        \caption{Ditinguidor D}

        $x \leftarrow \{0,1\}^n$

        $m_0 \leftarrow (x, x-1,x-2,\cdots,x-(l - 1))$ mod $2^n$

        $m_1 \leftarrow \{0\}^{nl}$

        $(IV, c_1, c_2,\cdots,c_l) \leftarrow O(m_0,m_1)$

        \If{todos $c_i$, com $ 1 \leq i \leq l$, forem iguais}{
            \Return 0
        }\Else{
            \Return 1
        }
    \end{algorithm}

    A vantagem pode ser calculada como segue.

    $Adv = |Pr(D^O = 0 | b = 0) - Pr(D^O = 0 |b = 1)|$

    $Pr(D^O = 0 | b = 0) = Pr(c_1 = c_2 = \cdots = c_l | b = 0) = 1$

    Como $m_1 = \{0\}^{nl}$, todo $c_i$, com $1 \leq i \leq l$, é $F_k(IV + i)$. Como $F_k$ é uma $PRP$ e toda entrada é diferente, o resultado de $F_k$ nas $i$ consultas será diferente, então $Pr(D^O = 0 |b = 1) = 0$.

    Logo,

    $Adv = |1 - 0| = 1$. Como $1 \notin negl$, $Adv \notin negl$.

\textbf{Questão 3}

    (Prova por contra-positiva) Se $\hat{\Pi}$ NÃO é CPA-Seguro, então $\Pi$ NÃO é CPA-Seguro.

    Desta forma, como $\hat{\Pi}$ NÃO é CPA-Seguro, existe um distinguidor $\hat{D}$ capaz de vencer o desafio CPA com vantagem não negligenciável dado $c$ e $t$. Agora construiremos um distinguidor $D$ para $\Pi$ da seguinte forma.

    \begin{algorithm}[H]
        \caption{Distinguidor D}

        Receba $m_0$ e $m_1$ de $\hat{D}$

        $c \leftarrow O((m_0, m_1))$

        $t \leftarrow g(c)$

        Retorne $(c,t)$ para $\hat{D}$

        Receba $b'$ de $\hat{D}$

        \Return $b'$
    \end{algorithm}

    A vantagem pode ser calculada como segue.

    $Adv^D = |Pr(D^O = 0 | b = 0) - Pr(D^O = 1 | b = 0)|$

    Como $Pr(D^O = 0 | b = 0) = Pr(\hat{D}^{\hat{O}} = 0 | b = 0)$ e $Pr(D^O = 1 | b = 0) = Pr(\hat{D}^{\hat{O}} = 1 | b = 0)$, temos que

    $Adv^D = |Pr(D^{\hat{O}} = 0 | b = 0) - Pr(D^{\hat{O}} = 1 | b = 0)| = Adv^{\hat{D}}$. Como $Adv^{\hat{D}} \notin negl$ (pois $\hat{D}$ é um distinguidor assumido na hipótese), então  $Adv^D \notin negl$ também. 


\textbf{Questão 4}

\textbf{Item A)}

    Seja $c_0 = r$ e $c_1 = (F_k(r)\oplus m)$. Uma função $Dec$ pode ser definida como segue.

    $Dec_k(c_0,c_1) = F_k(c_0) \oplus c_1 = F_k(r) \oplus F_k(r) \oplus m = m \in \{0,1\}^n$

    Como $\forall k \in K, \forall m \in M, Pr(Dec_k(Enc_k(m)) = m) = 1$, então $Dec$, como definido acima, garante corretude.

\textbf{Item B)}

    A segurança CPA de $\Pi$ se apoia no fato de $F$ ser uma $PRF$. Logo, se $F$ é uma $PRF$, então $\Pi$ é CPA-seguro.

    (Prova por contra-positiva) Se $\Pi$ NÃO é CPA-seguro, então $F$ NÃO é uma $PRF$.

    Como $\Pi$ não é CPA-seguro, então existe um distinguidor $D'$ capaz de vencer o desafio CPA com vantagem não negligenciável dado $(r, F_k(r) \oplus m)$. Desta forma, construiremos um distinguidor $D$, para mostrar que $F$ não é uma $PRF$, como segue. 

    \begin{algorithm}[H]
        \caption{Distinguidor D}

        Receba $m_0$ e $m_1$ de $D'$

        Seja $b$ um bit aleatório $\in \{0,1\}$

        $r \leftarrow \{0,1\}^n$

        $y \leftarrow O(r) \oplus m_b$

        Retorne $(r,y)$ para $D'$

        Receba $b'$ de $D'$

        \If{$b' = b$}{
            \Return 1
        }\Else{
            \Return 0
        }
    \end{algorithm}

    A vantagem pode ser calculada como segue.

    $Adv = |Pr(D^{F(.)} = 1 | O = F) - Pr(D^{f(.)} = 1 | O = f)|$

    $Pr(D^{F(.)} = 1 | O = F) = Pr(b' = b) = \frac{1}{2} + Adv^{D'}$ que $\notin negl$

    Como $f$ é uma função aleatória, $D'$ não sabe distinguir nesse cenário, logo $Pr(D^{f(.)} = 1 | O = f) = Pr(b' = b) = \frac{1}{2}$.

    Portanto,

    $Adv = |(\frac{1}{2} + Adv^{D'})  - \frac{1}{2}| = Adv^{D'} \notin negl$, então $F$ não é uma $PRF$.

\newpage

\textbf{Questão 5}

\textbf{Item A)}

    Para provarmos que $(G,\oplus)$ é um grupo, precisamos provar que há associatividade, o elemento neutro e que todo elemento em G tem um inverso, isto é:

    Associatividade: $\forall a,b,c \in G, (a \oplus b) \oplus c = a \oplus (b \oplus c)$

    Elemento neutro: $\exists e \in G \text{ tal que } \forall a \in G, a \oplus e = e \oplus a = a$

    Inverso: $\forall a \in G, \exists a^{-1} \in G, a \oplus a^{-1} = a^{-1} \oplus a = e$

    Associatividade: Seja $A$, $B$ e $C \in G$, ou seja, $A$, $B$ e $C \in \{0,1\}^n$. Como a operação $\oplus$ opera bit a bit, então temos que $(A \oplus B)_i = A_i \oplus B_i$, com $1 \leq i \leq n$. De forma análoga, temos que $((A \oplus B) \oplus C)_i = (A_i \oplus B_i) \oplus C_i$ e que $(A \oplus (B \oplus C))_i = A_i \oplus (B_i \oplus C_i)$. Como a operação $\oplus$ é associativa sobre bits ($\oplus$ sobre bits equivale a soma mod 2), temos que $(A_i \oplus B_i) \oplus C_i = A_i \oplus (B_i \oplus C_i)$. Como a igualdade é válida para todos os $i$ bits, temos que $(A \oplus B) \oplus C = A \oplus (B \oplus C)$.

    Elemento neutro: Como dito anteriormente, a operação $\oplus$ sobre bits equivale a soma mod 2. Pensando dessa forma, o elemento neutro $e = (0,0,\cdots, 0) \in G$. $\forall A \in G, A \oplus e = e \oplus A = A$, pois para todo bit $A_i$, com $1 \leq i \leq n$, $A_i \oplus 0 = A_i + 0 \mod 2 = A_i$ e $0 \oplus A_i = 0 + A_i \text{ mod } 2 = A_i$. Como a igualdade é valida para todos os $i$ bits, temos que $A \oplus e = A$ e $e \oplus A = A$, portanto $\forall A \in G, A \oplus e = e \oplus A = A$

    Inverso: Ainda se valendo da ideia de que a operação $\oplus$ sobre bits equivale a soma mod 2, o inverso de qualquer elemento $A \in G$ é ele mesmo, pois para cada $A_i$, com $1 \leq i \leq n$, temos que $A_i \oplus A_i = A_i + A_i \text{ mod } 2 = 0$. Observe que $A_i$ pode ser 0 ou 1, caso $A_i = 0$ temos $0 + 0 \text{ mod } 2 = 0$, caso $A_i = 1$ temos $1 + 1 \text{ mod } 2 = 0$, como as igualdades valem para todos os $i$ bits, então, para qualquer $A \in G, A^{-1} = A$, pois $A \oplus A = e$. 

    Portanto, $(G, \oplus)$ é um grupo.

\textbf{Item B)}

    Se a operação for trocada de $\oplus$ para o operador AND, $(G,\&)$ não é um grupo. Neste caso, o elemento neutro seria $e = (1,1,\cdots,1) \in \{0,1\}^n$. Para que qualquer $A \text{ \& } A^{-1} = A^{-1} \text{ \& } A = e$, com $A \in G$, isso somente seria verdade para strings da seguinte forma $A = (1, 1, \cdots, 1) \in \{0,1\}^n$, pois, se $A$ tiver qualquer bit $A_i = 0$, não importa qual outro bit esteja sendo operado, o resultado seria 0 e como $e = (1,1,\cdots, 1)$, $A \text{ \& } B \neq e$, para qualquer $B \in G$.
    
\textbf{Questão 6}

\textbf{Item A)}

    A composição de funções é definida como segue. Para qualquer par de funções $f$ e $g$, $f \circ g = f(g(x))$

    Associatividade: Seja $f$, $g$ e $h \in G$.

    $f \circ (g \circ h) = f(x) \circ g(h(x)) = f(g(h(x)))$

    $(f \circ g) \circ h = f(g(x)) \circ h(x) = f(g(h(x)))$

    Como $f \circ (g \circ h) = (f \circ g) \circ h$, então o conjunto $G$ é associativo.

    Elemento neutro: Seja $f$ e $e \in G$, tal que $e(x) = x$, isto é, $e$ é a função identidade.

    $f \circ e = f(e(x)) = f(x)$

    $e \circ f = e(f(x)) = f(x)$

    Como $f \circ e = e \circ f = f$, então o conjunto $G$ possui o elemento neutro.

    Inverso: Como $G$ é o conjunto de todas as funções $f$ tal que $f$ é bijetora, isso quer dizer que, para qualquer $f \in G$, existe $f^{-1} \in G$.

    $f \circ f^{-1} = f(f^{-1}(x)) = e$

    $f^{-1} \circ f = f^{-1}(f(x)) = e$

    Como as esquações acima são verdade para qualquer $f \in G$, então, para qualquer função, existe uma inversa tal que, quando composto, resulta no elemento neutro de $G$.

    Como $(G, \circ)$ é associativo, tem o elemento neutro e todo elemento tem um inverso correspondente, então $(G, \circ)$ é um grupo.

\textbf{Item B)}

    Como $|S| = 7$ e a operação é a composição de funções, temos que $G = S_7$, isto é, $G$ é o conjunto de todas as permutações possíveis de 7 elementos. Então agora a questão se resume em saber se existe $g \in G$ tal que $g^7 = id$ (ou seja, $g$ composto consigo mesmo 7 vezes tal que, na sétima vez, $g^7(n) = n$). A reposta é sim, pois como $G$ é o conjunto de todas as permutações possíveis de $S \rightarrow S$, então existe um $g \in G$ tal que $g(n) = n + 1 \mod 7$. Perceba que $g$ é um ciclo de tamanho 7 em todos os elementos de $S$. Portanto, $ord(g) = 7$. 

\textbf{Item C)}

    A ordem de $g$, como definido, é $ord(g) = n$, pois $g^t = i + t$ com $t < n$, ou seja, $g^t$ avança uma posição no ciclo e, quando $t = n$, $g^t = 1$, isto é, retorna ao início do ciclo.

\textbf{Item D)}
    
    Sabemos que $h = g^x$ e que $h(2) = 5$, então $g^x(2) = 5$. Como $g^x(i) = i + x$, temos que $g^x(2) = 2 + x$, o que implica que $2 + x = 5$, concluindo que $x = 3$. Portanto, $\log_g(h) = 3.$

\textbf{Questão 7}

\textbf{Item A)}

    Como $h_1 \text{ e } h_2 \in \langle g \rangle$, então $h_1 = g^x$ e $h_2 = g^y$.

    Com isso, temos que $\log_g(h_1) = \log_g(g^x) = x$ e $\log_g(h_2) = \log_g(g^y) = y$.

    Então, $\log_g(h_1) + \log_g(h_2) = \log_g(g^x) + \log_g(g^y) = x + y$ e 
    $\log_g(h_1 \cdot h_2) = \log_g(g^x \cdot g^y) = \log_g(g^{x+y}) = x + y$.

    Como $log_g(h_1) + \log_g(h_2) = x + y$ e $\log_g(h_1 \cdot h_2) = x + y$, temos que $\log_g(h_1 \cdot h_2) = \log_g(h_1) + \log_g(h_2)$

\textbf{Item B)}

    $\log_g(h_1^n) = \log_g((g^x)^n) = \log_g(g^{x \cdot n}) = x \cdot n = n \cdot x$

    $n \cdot \log_g(h_1) = n \cdot \log_g(g^x) = n \cdot x$

    Como $n \cdot \log_g(h_1) = n \cdot x$ e $\log_g(h_1^n) = n \cdot x$, então $n \cdot \log_g(h_1) = \log_g(h_1^n)$


\textbf{Questão 8}

    Se $b = 0$, então escolhe $y \in G$ uniformemente e computa $c_1 = g^y$ e $c_2 = h^y$. Texto cifrado $= (c_1,c_2)$

    Se $b = 1$, então escolhe $y, z \in G$ uniformemente e computa $c_1 = g^y$ e $c_2 = g^z$. Texto cifrado $= (c_1,c_2)$

\textbf{Item A)}

    Mostrar que é possível recuperar $b$ dado conhecimento sobre $sk = x$

    \begin{algorithm}[H]
        \caption{$Dec_k$}
        \KwIn{$(c_1,c_2)$}
        \KwOut{$b$}

        $y \leftarrow c_1^x$

        \If{$y = c_2$}{
            \Return 0
        }\Else{
            \Return 1
        }
        
    \end{algorithm}    

    Como $h = g^x$, temos que $c_2$ pode assumir 2 valores, sendo $(g^x)^y$ ou $(g^x)^z$. Quando a linha 1 calcula $c_1^x$, o resultado é $(g^y)^x = (g^x)^y$. Desta forma, caso $y = c_2$, isto é, $(g^y)^x = (g^x)^y$, significa dizer que $b = 0$, caso contrário, $b = 1$.

\textbf{Item B)}

    A questão basicamente pede para provar que se é DDH difícil sobre $G$, então é o esquema CPA-Seguro, desta forma, construirei a prova como segue.

    (Prova por contra-positiva) Se o esquema não é CPA-Seguro, então o DDH é facil sobre $G$. Para esta prova, assuma que a chave pública ($pk = (G,q,g,h)$ com $h = g^x$) é conhecida por todos, afinal, ela é pública.

    Se o esquema não é CPA-Seguro, então quer dizer que existe um algoritmo $PPT$ $D$ capaz de vencer o jogo CPA, isto é, $D$ é capaz de saber qual mensagem foi cifrada ao interagir com o oráculo.

    Seja $A$ o algoritmo capaz de resolver o DDH definido como segue.


    \begin{algorithm}[H]
        \caption{$A$}
        \KwIn{$(G, q, g, g^x, g^y, g^z)$}
        \KwOut{$0$ se $z = xy$ ou $1$ caso contrário}

        Receba $m_0$ e $m_1$ de $D$

        $b \leftarrow D(g^y, g^z)$

        \Return b        
    \end{algorithm}    

    A vantagem de $A$ pode ser calculada como segue.

    $Adv^A = |Pr(A = 0 | z=xy) - Pr(A = 0 | z \neq xy) |$

    Note que $Pr(A = 0 | z=xy) = Pr(D = 0 | b=0)$ e $Pr(A = 0 | z \neq xy) = Pr(D = 0 | b = 1)$. Desta forma temos que

    $Adv^A = |Pr(D = 0 | b=0) - Pr(D = 0 | b = 1)| = Adv^D$

    Como assumimos, por hipótese, que $D$ vence o jogo CPA com vantagem não negligenciável e $Adv^A = Adv^D$, então $Adv^A \notin negl$.

\textbf{Questão 9}

    Seja $(c_1, c_2)$ o texto cifrado de uma mensagem $m$. Como $c_1 = g^r$ e $c_2 = m \cdot h^r$ com $h = g^x$, podemos definir $(c_1', c_2')$ como segue.

    $c_1' = c_1$

    $c_2'= c_2 \cdot u = m \cdot h^r \cdot u$

    Note que, ao desencriptografar usando $Dec_k((c_1',c_2'))$, teremos

    $Dec_k((c_1',c_2')) = c_2' \cdot ((c_1')^x)^{-1} = m \cdot (g^x)^r \cdot u \cdot ((g^r)^x)^{-1} = m \cdot u$.

\textbf{Questão 10}

    Prove que se DDH é difícil sobre $G$, então CDH também é difícil sobre $G$.

    *Relembrando: 
    \begin{itemize}
        \item DLP $\rightarrow$ dado $g^x$, recuperar $x$
        \item CDH $\rightarrow$ dado $g^x$ e $g^y$, calcular $g^{xy}$
        \item DDH $\rightarrow$ dado $(G,q,g,g^x,g^y,g^z)$, descobrir se $z = xy$ ou se $z$ é um elemento aleatório em $\mathbb{Z}_q$
    \end{itemize}

    (Prova por contra-positiva) Se CDH é fácil sobre $G$, então DDH também é fácil sobre $G$.

    Se CDH é fácil, significa dizer que existe um algoritmo PPT $A$ que dado $g^x$ e $g^y$, retornará $g^{xy}$, portanto, seja $D$ o algoritmo PPT capaz de resolver o problema DDH definido como segue. Por conveniência, assuma que a chave pública $pk = (G,q,g,h)$, com $h = g^x$, é conhecida por todos os algoritmos, afinal, ela é pública.

    \begin{algorithm}[H]
        \caption{$D$}
        \KwIn{$(G, q, g, g^x, g^y, g^z)$}
        \KwOut{$1$ se $z = xy$ ou $0$ caso contrário}

        $t \leftarrow A(g^x, g^y)$

        \If{$t = g^z$}{
            \Return 1
        }\Else{
            \Return 0
        }    
    \end{algorithm}  

    Assumindo que $A$ sempre retorna $g^{xy}$, temos que a vantagem de $D$ pode ser calculada como segue.

    $Adv^D = |Pr(D = 1 | z = xy) - Pr(D = 1| z \leftarrow \mathbb{Z}_q)|$

    Note que $Pr(D = 1 | z = xy) = 1$, pois assumimos que $A$ tem êxito sempre, e $Pr(D = 1| z \leftarrow \mathbb{Z}_q) = \frac{1}{q}$, que é a probabilidade de haver colisão, isto é, a probabilidade de um $z$ aleatório ser igual a $xy$.

    Logo $Adv^D = |Pr(D = 1 | z = xy) - Pr(D = 1| z \leftarrow \mathbb{Z}_q)| = |1 - \frac{1}{q}|$. Como $q$ é a ordem de $g$ e geralmente é um número grande, $\frac{1}{q}$ é um número próximo a 0. Portanto, $Adv^D \approx 1$ que $\notin negl$.

\textbf{Questão 11}

    Se o melhor algoritmo para resolver o DLP em $G$ tem complexidade $2^{(\log q)\frac{1}{3}}$ e queremos ter 128 bits de segurança, façamos o seguinte

    $2^{(\log q)\frac{1}{3}} = 2^{128} \rightarrow (\log q)^\frac{1}{3} = 128 \rightarrow \sqrt[3]{\log q} = 128 \rightarrow \log q = (2^{7})^3 \rightarrow \log q = 2^{21}$.

    Portanto, a chave $sk$ precisa ter $2^{21}$ bits de tamanho.

    Como $h = g^x$ é um elemento no grupo, então $h$ tem tamanho igual $10 \cdot \log q \rightarrow 10 \cdot 2^{21}$.

    Os textos cifrados $(c_0,c_1)$ ambos são pertencentes ao grupos, logo, para qualquer mensagem $m$, o texto cifrado terá tamanho $2 \cdot 10 \cdot \log q \rightarrow 20 \cdot 2^{21}$
\end{document}